%----------------------------------------------------------------------------------------
% AFSTAND EN HOOFDASSEN. Alle opgaven zijn eigen werk; alle gegevens zijn verzonnen.
%
% Opbouw: Probleem, Start, leerstof (theorie, werkwijze of uitleg, voorbeeld, korte oefening), en op een
% nieuwe pagina de Opdrachten (door elkaar, stijgende moeilijkheid), Koppel en Test jezelf.
% Labels van dit hoofdstuk beginnen met h3-.
%----------------------------------------------------------------------------------------

\mychapter{3}{Afstand en hoofdassen}

\begin{probleem}
\naast{Van alle leerlingen van het zesde jaar kent de school de punten voor wiskunde ($W$) en fysica ($F$), in procent. De figuur toont een deel van de leerlingen. Voor het hele jaar is het gemiddelde punt $(62, 58)$ en de covariantiematrix
\[ S = \um{rr}{153 & 96 \\ 96 & 97} . \]
Lien haalt $83$ voor wiskunde en $80$ voor fysica, Ruben $74$ en $42$. De klassenraad wil weten welke leerling het meest afwijkt van wat gewoon is. Lien ligt het verst van het gemiddelde punt. Maar is haar combinatie ook het meest ongewoon? Hoe meet je een afstand tot het centrum van een wolk, als de variabelen elk anders spreiden en samenhangen?}{\begin{tikzpicture}[x=.055cm,y=.055cm,font=\scriptsize]
  \diagram[$W$][$F$]{30}{10}{95}{30}{10}{90}
  \xas{30,50,...,90}\yas{30,50,...,90}
  \punten[inkt]{(58,63) (49,45) (67,59) (55,53) (59,44) (43,44) (54,51) (65,63) (53,50) (62,62) (77,60) (59,52) (87,72) (71,57) (80,72) (52,49) (70,66) (54,50) (52,57) (67,57) (57,56) (53,43) (72,61) (82,74) (58,50) (43,49) (53,54) (55,44) (53,55) (82,75)}
  \gempunt{62}{58}
  \fill[fel!80!black] (83,80) circle (2.4pt) node[above left,text=black] {Lien};
  \fill[fel!80!black] (74,42) circle (2.4pt) node[below right,text=black] {Ruben};
\end{tikzpicture}}
\end{probleem}

\lkr{Leerplan Complementair (III-WiCo-d): LPD 53, met de wenken over de afstand van Mahalanobis (een maat voor de afstand tot het centrum die rekening houdt met de spreiding en de correlaties, en het verschil met de Euclidische afstand) en over eigenwaarden, eigenvectoren en diagonaliseren (de covariantiematrix is diagonaliseerbaar; met een basis van genormeerde eigenvectoren worden de variabelen ongecorreleerd, en wordt de afstand van Mahalanobis een Euclidische afstand). Eigenwaarden zijn zelf een module van Complementair (Matrices, LPD 30). Wie die module zag, kan sectie 2 snel doen; wie niet, krijgt hier wat nodig is, beperkt tot symmetrische $2 \times 2$-matrices, met de bewijzen. Voorkennis uit pakket 4: Matrices (inverse, determinant, homogene stelsels), Ruimtemeetkunde (scalair product, eenheidsvector, loodrechte stand) en hoofdstuk 2. De ellipsen verwijzen naar Kegelsneden, maar steunen er niet op. Extra: de hoofdcomponenten (de eerste hoofdas als richting met de grootste spreiding) en de variantie in elke richting. Met ICT: \texttt{numpy.linalg.eigh(S)} geeft de eigenwaarden van klein naar groot en de genormeerde eigenvectoren als kolommen; GeoGebra heeft \texttt{Eigenvalues} en \texttt{Eigenvectors}. Controleer bij ICT altijd het teken: een eigenvector mag je vermenigvuldigen met $-1$.}

Hoe ver ligt een punt van het centrum van een puntenwolk? Met een lat meet je de \textit{Euclidische} afstand, maar die houdt geen rekening met de vorm van de wolk. In een langwerpige wolk is een stap in de lengte iets gewoons, en een even grote stap dwars op de wolk uitzonderlijk. In dit hoofdstuk maak je een afstand die de covariantiematrix gebruikt: de \textit{afstand van Mahalanobis}. Om te begrijpen wat ze doet, zoek je de richtingen waarin de wolk het meest en het minst spreidt: de \textit{hoofdassen}. Dat zijn de eigenvectoren van de covariantiematrix. In het assenstelsel van de hoofdassen zijn de variabelen ongecorreleerd, en wordt de afstand van Mahalanobis een gewone afstand. Op het einde beslis je wie van Lien en Ruben het meest afwijkt.

\begin{oefeningen}[Start]
\oef{h3-start-1}

\needspace{12\baselineskip}
\oef{h3-verken}
\end{oefeningen}

\lkr{De Verken toont al het hoofdidee van sectie 2: de spreiding in de richting $(2, 1)$ is de grootste, en $S$ vermenigvuldigt die vector met 6. Laat leerlingen raden welke richting de kleinste spreiding geeft (loodrecht erop).}

%========================================================================================
\section{Afstand tot het centrum}

\subsection{Waarom de lat niet volstaat}

Neem een wolk met centrum $\overline{\vec x}$. De Euclidische afstand van een punt $\vec x$ tot het centrum is de lengte van $\vec x - \overline{\vec x}$: bij twee variabelen $\sqrt{(x_1 - \xbar_1)^2 + (x_2 - \xbar_2)^2}$. Ze heeft twee gebreken.

Ten eerste hangt ze af van de eenheden. Meet je de lengte van een persoon in cm en zijn massa in kg, dan telt een verschil van 10 cm even zwaar als 10 kg. In mm zou de lengte bijna alles bepalen. Dat los je op met z-scores: deel elke afwijking door de standaardafwijking van die variabele. Zo krijg je de \textit{gestandaardiseerde} afstand $\sqrt{z_1^2 + z_2^2}$.

Ten tweede houdt ook die geen rekening met de samenhang. Hangen de variabelen positief samen, dan ligt de wolk schuin. Een punt in de lengterichting van de wolk is dan minder ongewoon dan een punt even ver ernaast. Een cirkel rond het centrum past niet bij de vorm van de wolk: je hebt een ellips nodig.

\figuren{\begin{tikzpicture}[x=.5cm,y=.5cm,font=\scriptsize]
  \draw[black!14,very thin] (-5,-4) grid (5,4);
  \draw[->,grijs] (-5,0) -- (5.4,0) node[right,text=black] {$x_1 - \xbar_1$};
  \draw[->,grijs] (0,-4) -- (0,4.4) node[above,text=black] {$x_2 - \xbar_2$};
  \draw[inkt,dashed] (0,0) circle (3);
  \ellips{0}{0}{4.899}{2}{26.565}
  \punten[inkt]{(-0.71,0.37) (-1.17,-0.82) (0.75,-0.79) (1.37,0.45) (1.11,-0.18) (2.57,1.93) (-0.53,0.29) (2.99,3.55) (-3.95,-1.27) (1.05,0.29) (-2.56,-0.08) (-3.18,-1.22) (3.09,-0.37) (-0.83,-2.15)}
  \fill[fel!80!black] (3,0) circle (2.6pt) node[below right,text=black] {$A$};
  \fill[fel!80!black] (0,3) circle (2.6pt) node[above left,text=black] {$B$};
  \node[inkt,anchor=west] at (0.8,-3.6) {cirkel: Euclidische afstand 3};
  \node[accent,anchor=west] at (-5.6,-3.6) {ellips: $\dM = 2$};
\end{tikzpicture}}

In de figuur liggen $A$ en $B$ allebei op Euclidische afstand 3 van het centrum. Maar $A$ ligt in de richting waarin de wolk sterk spreidt, $B$ in een richting waarin ze weinig spreidt. $B$ is dus het meest uitzonderlijk. De ellips volgt de vorm van de wolk: alle punten op de ellips zijn \textit{even ongewoon}. Welke afstand hoort bij zo'n ellips?

\subsection{De afstand van Mahalanobis}

Bij één variabele meet de z-score de afstand tot het gemiddelde in standaardafwijkingen:
\[ z^2 = \frac{(x - \xbar)^2}{s^2} = (x - \xbar)\,(s^2)^{-1}\,(x - \xbar) . \]
Bij meer variabelen vervang je de variantie $s^2$ door de covariantiematrix $S$, en het delen door de inverse matrix.

\begin{theorie}
Is $S$ inverteerbaar, dan is de \textbf{\textcolor{accent}{afstand van Mahalanobis}} van een punt $\vec x$ tot het centrum $\overline{\vec x}$, beide als kolom,

\vspace{1mm}
\qquad $\formule{\dM(\vec x) = \sqrt{(\vec x - \overline{\vec x})\tr\,S^{-1}\,(\vec x - \overline{\vec x})}}$
\end{theorie}

De Indiase statisticus Prasanta Chandra Mahalanobis voerde ze in 1936 in, om groepen te vergelijken op basis van veel lichaamsmaten tegelijk. Onder de wortel staat nooit een negatief getal: dat volgt in sectie 3. Bij twee variabelen is $S^{-1} = \frac{1}{\det S}\km{s_2^2 & -s_{12} \\ -s_{12} & s_1^2}$, en met de afwijkingen $d_1 = x_1 - \xbar_1$ en $d_2 = x_2 - \xbar_2$ wordt de formule
\[ \dM^2 = \frac{s_2^2\,d_1^2 - 2s_{12}\,d_1d_2 + s_1^2\,d_2^2}{s_1^2s_2^2 - s_{12}^2} . \]

\voorbeeld{}{Een wolk heeft centrum $(0, 0)$ en covariantiematrix $S = \km{5 & 2 \\ 2 & 2}$. Bereken de Euclidische afstand en de afstand van Mahalanobis voor $A = (3, 0)$ en $B = (0, 3)$ uit de figuur.\label{vb:h3-mahalanobis}}{%
\uitleg{We berekenen eerst $S^{-1}$ met de formule uit Matrices, en dan voor elk punt een rij maal een matrix maal een kolom.}
\begin{stappenbreed}
S^{-1} = \tfrac16\um{rr}{2 & -2 \\ -2 & 5} & $\det S = 10 - 4 = 6$. \\[3.5ex]
A\colon\ \dM^2 = \tfrac16\,(3,\ 0)\um{rr}{2 & -2 \\ -2 & 5}\um{c}{3 \\ 0} = \tfrac{18}{6} = 3, \quad \dM \approx 1{,}73 & Alleen het element linksboven telt. \\[3.5ex]
B\colon\ \dM^2 = \tfrac16 \cdot 5 \cdot 9 = 7{,}5, \quad \dM \approx 2{,}74 & Alleen het element rechtsonder telt.
\end{stappenbreed}
De Euclidische afstand is voor beide 3, maar $B$ ligt veel verder in de zin van Mahalanobis. De wolk spreidt in de $x_2$-richting minder ($s_2^2 = 2$) dan in de $x_1$-richting ($s_1^2 = 5$).}

\begin{werkwijze}
Bereken de afwijking $\vec d = \vec x - \overline{\vec x}$ en de inverse $S^{-1}$. Bereken $S^{-1}\vec d$, dan het scalair product met $\vec d$, en neem de wortel. Bij twee variabelen kan ook de formule met $\det S$. Een afstand van Mahalanobis meet je in \textit{standaardafwijkingen}: een waarde boven 3 is bij een normaal verdeelde wolk zeldzaam.
\end{werkwijze}

\begin{oefeningen}
\oef{h3-mahalanobis}
\end{oefeningen}

\subsection{Twee eigenschappen}

Waarom is dit een goede afstand? Twee eigenschappen tonen dat ze de gebreken van de Euclidische afstand oplost.

\begin{theorie}
\begin{stelling}{Ongecorreleerde variabelen}
Is $S = \diag(s_1^2, \ldots, s_p^2)$, dan is \qquad $\formule{\dM = \sqrt{z_1^2 + \cdots + z_p^2}}$ \qquad met $z_j = \frac{x_j - \xbar_j}{s_j}$
\end{stelling}

\begin{bewijs}
De inverse van een diagonaalmatrix is $S^{-1} = \diag\bigl(\frac{1}{s_1^2}, \ldots, \frac{1}{s_p^2}\bigr)$, dus
\bewijsregel{\dM^2 = \sum_{j=1}^{p} \frac{(x_j - \xbar_j)^2}{s_j^2} = \sum_{j=1}^{p} z_j^2 .}
\end{bewijs}

\begin{stelling}{Onafhankelijk van de eenheden}
Ontstaan nieuwe variabelen met $\vec y = A\vec x + \vec b$ en een inverteerbare matrix $A$, dan heeft elk punt dezelfde afstand van Mahalanobis in de nieuwe variabelen als in de oude.
\end{stelling}

\begin{bewijs}
De afwijking wordt $\vec y - \overline{\vec y} = A\vec d$, en de covariantiematrix $ASA\tr$ (hoofdstuk 2). Omdat $A\tr(A^{-1})\tr = (A^{-1}A)\tr = I$, is $(A\tr)^{-1} = (A^{-1})\tr$, en dus $(ASA\tr)^{-1} = (A^{-1})\tr S^{-1}A^{-1}$. Dan is
\bewijsregel{(A\vec d)\tr (ASA\tr)^{-1}(A\vec d) = \vec d\tr A\tr (A^{-1})\tr S^{-1} A^{-1}A\,\vec d = \vec d\tr S^{-1}\vec d .}
\end{bewijs}
\end{theorie}

De eerste stelling zegt: zonder samenhang is de afstand van Mahalanobis de gestandaardiseerde afstand. Met samenhang verschilt ze. De tweede stelling zegt meer dan alleen \textit{onafhankelijk van de eenheden} (dan is $A$ een diagonaalmatrix): ook als je overstapt op het totaal en het verschil, of op een ander assenstelsel, blijft de afstand dezelfde. Dat gebruiken we in sectie 3.

\begin{oefeningen}
\oef{h3-eigenschappen}
\end{oefeningen}

\subsection{Ellipsen van gelijke afstand}

\naast{\begin{minipage}{9.2cm}
Alle punten met $\dM = k$ voldoen bij twee variabelen aan $s_2^2d_1^2 - 2s_{12}d_1d_2 + s_1^2d_2^2 = k^2\det S$. Dat is een vergelijking van de tweede graad in $d_1$ en $d_2$, en de figuur toont dat het een ellips is rond het centrum. Rechts zie je de ellipsen met $\dM = 1$, $2$ en $3$ voor de wolk van het voorbeeld. Hoe groter $k$, hoe groter de ellips, maar de vorm en de stand blijven dezelfde. Waarom liggen de assen van de ellipsen schuin, en hoe lang zijn ze? Dat zoeken we in de volgende sectie uit.
\end{minipage}}{\begin{tikzpicture}[x=.4cm,y=.4cm,font=\scriptsize]
  \draw[black!14,very thin] (-7,-5) grid (7,5);
  \draw[->,grijs] (-7,0) -- (7.4,0) node[right,text=black] {$d_1$};
  \draw[->,grijs] (0,-5) -- (0,5.4) node[above,text=black] {$d_2$};
  \ellips{0}{0}{2.449}{1}{26.565}
  \ellips{0}{0}{4.899}{2}{26.565}
  \ellips{0}{0}{7.348}{3}{26.565}
  \punten[inkt]{(-0.71,0.37) (-1.17,-0.82) (0.75,-0.79) (1.37,0.45) (1.11,-0.18) (2.57,1.93) (-0.53,0.29) (2.99,3.55) (-3.95,-1.27) (1.05,0.29) (-2.56,-0.08) (-3.18,-1.22) (3.09,-0.37) (-0.83,-2.15)}
  \node[accent,fill=white,inner sep=1pt] at (2.1,1.45) {1};
  \node[accent,fill=white,inner sep=1pt] at (3.6,3.0) {2};
  \node[accent,fill=white,inner sep=1pt] at (5.1,4.45) {3};
\end{tikzpicture}}

%========================================================================================
\section{Eigenvectoren van de covariantiematrix}

\subsection{De richting met de grootste spreiding}

In hoofdstuk 2 zag je dat $\vec a\tr S\vec a$ de spreiding van de wolk in de richting van een eenheidsvector $\vec a$ meet. Neem $\vec a = (\cos\theta, \sin\theta)$ en laat $\theta$ draaien. Voor $S = \km{5 & 2 \\ 2 & 2}$ is
\[ V(\theta) = \vec a\tr S\vec a = 5\cos^2\theta + 4\sin\theta\cos\theta + 2\sin^2\theta . \]

\naast{\begin{minipage}{8.2cm}
De grafiek toont $V(\theta)$ van $0^\circ$ tot $180^\circ$ (na $180^\circ$ herhaalt alles, want $-\vec a$ geeft dezelfde richting). Het maximum is 6, bij $\theta \approx 26{,}6^\circ$: de richting $(2, 1)$ van de Verken. Het minimum is 1, precies $90^\circ$ verder, in de richting $(-1, 2)$. Dat zijn de richtingen van de assen van de ellipsen. En in de Verken zag je iets merkwaardigs: $S$ maakt van de vector $(2, 1)$ een veelvoud van zichzelf,
\[ S\km{2 \\ 1} = \km{12 \\ 6} = 6\km{2 \\ 1} . \]
Het getal 6 is de grootste spreiding zelf. Toeval? Nee: de richtingen die een matrix niet verandert, zijn de sleutel tot dit hoofdstuk.
\end{minipage}}{\begin{tikzpicture}[x=.03cm,y=.42cm,font=\scriptsize]
  \draw[black!14,very thin,xstep=30,ystep=1] (0,0) grid (180,7);
  \draw[->,grijs] (0,0) -- (190,0) node[right,text=black] {$\theta$};
  \draw[->,grijs] (0,0) -- (0,7.5) node[above,text=black] {$V$};
  \foreach \t in {30,90,150}{\node[below] at (\t,0) {$\t^\circ$};}
  \foreach \t in {1,2,...,6}{\node[left] at (0,\t) {$\t$};}
  \draw[accent,thick] plot[domain=0:180,samples=90,smooth] (\x,{3.5+1.5*cos(2*\x)+2*sin(2*\x)});
  \draw[hulplijn] (26.57,0) -- (26.57,6);
  \draw[hulplijn] (116.57,0) -- (116.57,1);
  \fill[fel!80!black] (26.57,6) circle (2pt);
  \fill[fel!80!black] (116.57,1) circle (2pt);
\end{tikzpicture}}

\subsection{Eigenwaarden en eigenvectoren}

\begin{theorie}
Een vector $\vec v \neq \vec 0$ is een \textbf{\textcolor{accent}{eigenvector}} van een vierkante matrix $S$ als er een getal $\lambda$ is met

\vspace{1mm}
\qquad $\formule{S\vec v = \lambda\vec v}$

\vspace{1mm}
Het getal $\lambda$ is dan een \textbf{\textcolor{accent}{eigenwaarde}} van $S$.

\begin{stelling}{Karakteristieke vergelijking}
$\lambda$ is een eigenwaarde van $S$ als en slechts als \qquad $\formule{\det(S - \lambda I) = 0}$
\end{stelling}

\begin{bewijs}
$S\vec v = \lambda\vec v$ is hetzelfde als $(S - \lambda I)\vec v = \vec 0$, een homogeen stelsel. Het heeft een oplossing $\vec v \neq \vec 0$ als en slechts als de matrix $S - \lambda I$ niet inverteerbaar is (Matrices), dus als en slechts als $\det(S - \lambda I) = 0$.
\end{bewijs}
\end{theorie}

\naast{\begin{minipage}{10.2cm}
Een eigenvector is dus een richting die de matrix niet draait: $S\vec v$ ligt op dezelfde rechte als $\vec v$. In de figuur is $\vec v = (2, 1)$ een eigenvector, met $S\vec v = (12, 6)$ in het verlengde. De vector $\vec w = (1, 0)$ is er geen: $S\vec w = (5, 2)$ draait weg. Is $\vec v$ een eigenvector, dan ook elk veelvoud $c\vec v$ met $c \neq 0$: we zoeken dus richtingen, en kiezen daarin een handige vector. Bij een $2 \times 2$-matrix is $\det(S - \lambda I)$ een veelterm van de tweede graad in $\lambda$, de \textit{karakteristieke veelterm}. Er zijn dus hoogstens twee eigenwaarden.
\end{minipage}}{\begin{tikzpicture}[x=.36cm,y=.36cm,font=\scriptsize]
  \draw[black!14,very thin] (0,0) grid (12,7);
  \draw[->,grijs] (0,0) -- (12.5,0);
  \draw[->,grijs] (0,0) -- (0,7.5);
  \draw[accent!50,dashed] (0,0) -- (13,6.5);
  \draw[->,accent,very thick] (0,0) -- (2,1) node[above left] {$\vec v$};
  \draw[->,accent,thick] (0,0) -- (12,6) node[above] {$S\vec v$};
  \draw[->,inkt,very thick] (0,0) -- (1,0) node[below] {$\vec w$};
  \draw[->,inkt,thick] (0,0) -- (5,2) node[below right] {$S\vec w$};
\end{tikzpicture}}

\voorbeeld{}{Bereken de eigenwaarden en eigenvectoren van $S = \km{5 & 2 \\ 2 & 2}$.\label{vb:h3-eigen}}{%
\uitleg{We lossen eerst de karakteristieke vergelijking op. Daarna zoeken we bij elke eigenwaarde een oplossing van het homogene stelsel $(S - \lambda I)\vec v = \vec 0$.}
\begin{stappenbreed}
\det\um{cc}{5 - \lambda & 2 \\ 2 & 2 - \lambda} = (5 - \lambda)(2 - \lambda) - 4 = \lambda^2 - 7\lambda + 6 & De karakteristieke veelterm. \\[3ex]
\lambda^2 - 7\lambda + 6 = 0 \iff \lambda = 6 \ \text{of}\ \lambda = 1 & De eigenwaarden. \\[1ex]
\lambda = 6\colon\ \um{rr}{-1 & 2 \\ 2 & -4}\vec v = \vec 0 \iff v_1 = 2v_2, \quad \vec v_1 = \um{c}{2 \\ 1} & Twee keer dezelfde vergelijking. \\[3.5ex]
\lambda = 1\colon\ \um{rr}{4 & 2 \\ 2 & 1}\vec v = \vec 0 \iff v_2 = -2v_1, \quad \vec v_2 = \um{r}{-1 \\ 2} & Kies $v_1 = -1$.
\end{stappenbreed}
Controle: $S\vec v_2 = (-5 + 4, -2 + 4) = (-1, 2)$. De eigenwaarden 6 en 1 zijn het maximum en het minimum van $V(\theta)$, en de eigenvectoren staan loodrecht op elkaar: $\vec v_1 \cdot \vec v_2 = -2 + 2 = 0$.}

\begin{werkwijze}
Schrijf $\det(S - \lambda I)$ uit en los de vergelijking op: dat geeft de eigenwaarden. Vul elke eigenwaarde in $(S - \lambda I)\vec v = \vec 0$ in. De twee vergelijkingen zijn dan veelvouden van elkaar: kies een oplossing met eenvoudige getallen. Controleer met $S\vec v = \lambda\vec v$.
\end{werkwijze}

\begin{oefeningen}
\oef{h3-eigen}
\end{oefeningen}

\subsection{Symmetrische matrices}

In het voorbeeld waren de eigenwaarden reëel en positief, en stonden de eigenvectoren loodrecht op elkaar. Bij een willekeurige matrix hoeft dat niet: $\km{0 & -1 \\ 1 & 0}$ draait elke vector over $90^\circ$ en heeft geen enkele reële eigenvector. Maar een covariantiematrix is symmetrisch, en dan gaat het altijd goed. Drie eigenschappen.

\begin{theorie}
\begin{stelling}{Reële eigenwaarden}
Een symmetrische $2 \times 2$-matrix heeft reële eigenwaarden.
\end{stelling}

\begin{bewijs}
Neem $S = \km{a & b \\ b & c}$. De karakteristieke vergelijking is $\lambda^2 - (a + c)\lambda + ac - b^2 = 0$, met discriminant
\bewijsregel{(a + c)^2 - 4(ac - b^2) = (a - c)^2 + 4b^2 \geqslant 0 .}
Ze is 0 alleen als $a = c$ en $b = 0$: dan is $S = aI$, met één eigenwaarde $a$.
\end{bewijs}

\begin{stelling}{Loodrechte eigenvectoren}
Zijn $\vec v_1$ en $\vec v_2$ eigenvectoren van een symmetrische matrix $S$ bij verschillende eigenwaarden $\lambda_1 \neq \lambda_2$, dan is \qquad $\formule{\vec v_1 \cdot \vec v_2 = 0}$
\end{stelling}

\begin{bewijs}
Schrijf het scalair product als $\vec v_1\tr\vec v_2$. Omdat $S\tr = S$, is
\bewijsregel{\lambda_1\,\vec v_1\tr\vec v_2 = (S\vec v_1)\tr\vec v_2 = \vec v_1\tr S\tr\vec v_2 = \vec v_1\tr(S\vec v_2) = \lambda_2\,\vec v_1\tr\vec v_2 .}
Dus $(\lambda_1 - \lambda_2)\,\vec v_1\tr\vec v_2 = 0$, en omdat $\lambda_1 \neq \lambda_2$, is $\vec v_1\tr\vec v_2 = 0$.
\end{bewijs}

\begin{stelling}{Eigenwaarden van een covariantiematrix}
Elke eigenwaarde van een covariantiematrix is $\formule{\lambda \geqslant 0}$.
\end{stelling}

\begin{bewijs}
Uit $S\vec v = \lambda\vec v$ volgt $\vec v\tr S\vec v = \lambda\,\vec v\tr\vec v = \lambda\abs{\vec v}^2$. Links staat een variantie, die niet negatief is (hoofdstuk 2), en $\abs{\vec v}^2 > 0$.
\end{bewijs}
\end{theorie}

Het bewijs van de tweede eigenschap werkt voor elke grootte van $S$; de eerste geldt ook voor grotere symmetrische matrices, maar het bewijs valt buiten het leerplan. Bij een covariantiematrix is een eigenwaarde 0 als en slechts als $S$ niet inverteerbaar is: dan is er een richting zonder spreiding, zoals in hoofdstuk 2.

\begin{oefeningen}
\oef{h3-symmetrisch}
\end{oefeningen}

\subsection{Diagonaliseren}

Deel elke eigenvector door zijn norm: dan krijg je \textbf{genormeerde} eigenvectoren, eenheidsvectoren die loodrecht op elkaar staan. Zet ze als kolommen in een matrix $P$. Wat zegt dat over $S$?

\begin{theorie}
\begin{stelling}{De covariantiematrix is diagonaliseerbaar}
Zijn $\vec e_1$ en $\vec e_2$ genormeerde eigenvectoren van een symmetrische $2 \times 2$-matrix $S$ bij $\lambda_1$ en $\lambda_2$, loodrecht op elkaar, en is $P = \bigl[\,\vec e_1 \ \ \vec e_2\,\bigr]$ de matrix met die kolommen, dan is

\qquad $\formule{P\tr P = I}$ \qquad en \qquad $\formule{S = P\,D\,P\tr}$ \qquad met $D = \diag(\lambda_1, \lambda_2)$
\end{stelling}

\begin{bewijs}
Het element $(j, k)$ van $P\tr P$ is $\vec e_j \cdot \vec e_k$: 1 als $j = k$ (eenheidsvectoren) en 0 anders (loodrecht). Dus $P\tr P = I$, en $P^{-1} = P\tr$. De kolommen van $SP$ zijn $S\vec e_1 = \lambda_1\vec e_1$ en $S\vec e_2 = \lambda_2\vec e_2$, en dat zijn ook de kolommen van $PD$. Uit $SP = PD$ volgt $S = PDP^{-1} = PDP\tr$.
\end{bewijs}
\end{theorie}

Zijn de eigenwaarden gelijk, dan is $S = \lambda I$ en is elke vector een eigenvector: kies dan gewoon $P = I$. Ook bij $p$ variabelen bestaat er altijd zo'n matrix $P$ met loodrechte genormeerde eigenvectoren. Dat is de \textit{spectraalstelling}; het bewijs valt buiten het leerplan. Een matrix $P$ met $P\tr P = I$ heet \textbf{orthogonaal}: ze verandert geen lengtes en hoeken, want $(P\vec a) \cdot (P\vec b) = \vec a\tr P\tr P\,\vec b = \vec a \cdot \vec b$. Bij twee variabelen is ze een draaiing, of een draaiing met een spiegeling.

\voorbeeld{}{Diagonaliseer $S = \km{5 & 2 \\ 2 & 2}$ uit voorbeeld \ref{vb:h3-eigen}.\label{vb:h3-diagonaliseren}}{%
\uitleg{De eigenvectoren $(2, 1)$ en $(-1, 2)$ hebben allebei norm $\sqrt5$. We delen erdoor en zetten ze als kolommen in $P$.}
\begin{stappenbreed}
P = \tfrac{1}{\sqrt5}\um{rr}{2 & -1 \\ 1 & 2}, \quad D = \um{cc}{6 & 0 \\ 0 & 1} & Eerste kolom bij $\lambda_1 = 6$. \\[3.5ex]
PD = \tfrac{1}{\sqrt5}\um{rr}{12 & -1 \\ 6 & 2} & Elke kolom maal haar eigenwaarde. \\[3.5ex]
PDP\tr = \tfrac15\um{rr}{12 & -1 \\ 6 & 2}\um{rr}{2 & 1 \\ -1 & 2} = \tfrac15\um{rr}{25 & 10 \\ 10 & 10} = S & Controle: het klopt.
\end{stappenbreed}
De matrix $P$ is een draaiing over $26{,}6^\circ$: ze brengt de $x$-as naar de richting van $\vec e_1$.}

\begin{oefeningen}
\oef{h3-diagonaliseren}
\end{oefeningen}

%========================================================================================
\section{Een nieuw assenstelsel}

\subsection{Ongecorreleerde variabelen}

Neem de hoofdassen van de wolk als nieuwe assen: door het centrum, in de richting van $\vec e_1$ en $\vec e_2$. De nieuwe coördinaten van een punt zijn de projecties van zijn afwijking op die assen: $u_1 = \vec e_1 \cdot (\vec x - \overline{\vec x})$ en $u_2 = \vec e_2 \cdot (\vec x - \overline{\vec x})$. In matrixvorm: $\vec u = P\tr(\vec x - \overline{\vec x})$. Wat is de covariantiematrix van de nieuwe variabelen?

\begin{theorie}
\begin{stelling}{De hoofdassen maken de variabelen ongecorreleerd}
De variabelen $\vec u = P\tr(\vec x - \overline{\vec x})$ hebben als covariantiematrix \qquad $\formule{S_u = D = \diag(\lambda_1, \lambda_2)}$
\end{stelling}

\begin{bewijs}
Dit is een lineaire transformatie met $A = P\tr$, dus $S_u = P\tr S\,(P\tr)\tr = P\tr S P$ (hoofdstuk 2). Met $S = PDP\tr$ en $P\tr P = I$ is
\bewijsregel{S_u = P\tr P\,D\,P\tr P = D .}
\end{bewijs}
\end{theorie}

In het nieuwe assenstelsel zijn de variabelen dus \textit{ongecorreleerd}: de covariantie is 0. Hun varianties zijn de eigenwaarden. De nieuwe variabelen heten de \textbf{hoofdcomponenten} van de gegevens. Ze zijn lineaire combinaties van de oude: bij $S = \km{5 & 2 \\ 2 & 2}$ is $u_1 = \frac{1}{\sqrt5}(2d_1 + d_2)$ en $u_2 = \frac{1}{\sqrt5}(-d_1 + 2d_2)$.

\begin{theorie}
\begin{stelling}{De eerste hoofdas}
Voor elke eenheidsvector $\vec a$ is \qquad $\formule{\lambda_2 \leqslant \vec a\tr S\vec a \leqslant \lambda_1}$ \qquad met $\lambda_1 \geqslant \lambda_2$ de eigenwaarden. Het maximum wordt bereikt bij $\vec a = \vec e_1$.
\end{stelling}

\begin{bewijs}
Stel $\vec b = P\tr\vec a$. Dan is $\abs{\vec b} = \abs{\vec a} = 1$, want $P\tr$ is orthogonaal ($P\,P\tr = P\,P^{-1} = I$). Uit $S = PDP\tr$ volgt
\bewijsregel{\vec a\tr S\vec a = \vec b\tr D\,\vec b = \lambda_1b_1^2 + \lambda_2b_2^2 \leqslant \lambda_1(b_1^2 + b_2^2) = \lambda_1 ,}
en net zo $\geqslant \lambda_2(b_1^2 + b_2^2) = \lambda_2$. Bij $\vec a = \vec e_1$ is $\vec a\tr S\vec a = \lambda_1\vec e_1 \cdot \vec e_1 = \lambda_1$.
\end{bewijs}
\end{theorie}

Dat verklaart de grafiek van $V(\theta)$: de grootste spreiding is de grootste eigenwaarde, in de richting van de eerste eigenvector. Die richting heet de \textbf{eerste hoofdas}. Het deel $\frac{\lambda_1}{\lambda_1 + \lambda_2}$ van de totale variantie ligt langs die as. Bij $S = \km{5 & 2 \\ 2 & 2}$ is dat $\frac67 \approx 86~\%$: wie de twee variabelen wil samenvatten in één getal, neemt $u_1$, en verliest maar $14~\%$ van de spreiding. Dat is het idee van de \textit{hoofdcomponentenanalyse}, een van de meest gebruikte technieken voor gegevens met veel variabelen.

\voorbeeld{}{Een punt heeft de afwijking $\vec d = (3, 4)$ tot het centrum van de wolk met $S = \km{5 & 2 \\ 2 & 2}$. Bereken de coördinaten in het assenstelsel van de hoofdassen.\label{vb:h3-assen}}{%
\uitleg{De nieuwe coördinaten zijn de scalaire producten met de genormeerde eigenvectoren uit voorbeeld \ref{vb:h3-diagonaliseren}.}
\begin{stappenbreed}
u_1 = \tfrac{1}{\sqrt5}(2 \cdot 3 + 1 \cdot 4) = \tfrac{10}{\sqrt5} = 2\sqrt5 \approx 4{,}47 & Langs de eerste hoofdas. \\[1ex]
u_2 = \tfrac{1}{\sqrt5}(-1 \cdot 3 + 2 \cdot 4) = \tfrac{5}{\sqrt5} = \sqrt5 \approx 2{,}24 & Langs de tweede hoofdas. \\[1ex]
u_1^2 + u_2^2 = 20 + 5 = 25 = 3^2 + 4^2 & De Euclidische afstand blijft 5.
\end{stappenbreed}
Het punt ligt $4{,}47$ ver langs de eerste hoofdas en $2{,}24$ langs de tweede. Een draaiing verandert de afstand tot het centrum niet.}

\begin{oefeningen}
\oef{h3-assen}
\end{oefeningen}

\subsection{De afstand van Mahalanobis als gewone afstand}

Nu het verband met de eerste sectie. In de nieuwe coördinaten is de covariantiematrix de diagonaalmatrix $D$, en de afstand van Mahalanobis verandert niet bij een lineaire transformatie. Dan zegt de stelling over ongecorreleerde variabelen meteen wat de afstand is.

\begin{theorie}
\begin{stelling}{Mahalanobis in de hoofdassen}
Met $\vec u = P\tr(\vec x - \overline{\vec x})$ is \qquad $\formule{\dM^2 = \frac{u_1^2}{\lambda_1} + \frac{u_2^2}{\lambda_2}}$

\smallskip
Met $w_j = \frac{u_j}{\sqrt{\lambda_j}}$ is $\dM = \sqrt{w_1^2 + w_2^2}$: de Euclidische afstand in de coördinaten $(w_1, w_2)$.
\end{stelling}

\begin{bewijs}
De afstand blijft dezelfde in de variabelen $\vec u$ (de stelling over de eenheden, met $A = P\tr$). Hun covariantiematrix is $D$, een diagonaalmatrix, dus door de stelling over ongecorreleerde variabelen is $\dM^2 = \frac{u_1^2}{\lambda_1} + \frac{u_2^2}{\lambda_2}$, en dat is $w_1^2 + w_2^2$.
\end{bewijs}
\end{theorie}

De drie figuren tonen het in twee stappen. Links de wolk met haar hoofdassen. In het midden is ze gedraaid: de hoofdassen zijn de nieuwe assen en de wolk ligt recht. Rechts is elke as herschaald met $\sqrt{\lambda_j}$: de wolk is rond, en de ellipsen van gelijke afstand zijn cirkels geworden. Zo volgt ook de vorm van de ellipsen: $\dM = k$ betekent $\frac{u_1^2}{k^2\lambda_1} + \frac{u_2^2}{k^2\lambda_2} = 1$, een ellips met halve assen $k\sqrt{\lambda_1}$ langs $\vec e_1$ en $k\sqrt{\lambda_2}$ langs $\vec e_2$.

\figuren{\begin{tikzpicture}[x=.4cm,y=.4cm,font=\scriptsize]
  \draw[->,grijs] (-5.2,0) -- (5.2,0) node[right,text=black] {$d_1$};
  \draw[->,grijs] (0,-4) -- (0,4.4) node[above,text=black] {$d_2$};
  \ellips{0}{0}{2.449}{1}{26.565}
  \ellips{0}{0}{4.899}{2}{26.565}
  \draw[inkt,thick,->] (0,0) -- (4.47,2.24) node[right] {$\vec e_1$};
  \draw[inkt,thick,->] (0,0) -- (-1.12,2.24) node[above] {$\vec e_2$};
  \punten[inkt]{(-0.71,0.37) (-1.17,-0.82) (0.75,-0.79) (1.37,0.45) (1.11,-0.18) (2.57,1.93) (-0.53,0.29) (2.99,3.55) (-3.95,-1.27) (1.05,0.29) (-2.56,-0.08) (-3.18,-1.22) (3.09,-0.37) (-0.83,-2.15)}
  \node[below] at (0,-4.3) {oude assen};
\end{tikzpicture}\quad
\begin{tikzpicture}[x=.4cm,y=.4cm,font=\scriptsize]
  \draw[->,grijs] (-5.2,0) -- (5.2,0) node[right,text=black] {$u_1$};
  \draw[->,grijs] (0,-4) -- (0,4.4) node[above,text=black] {$u_2$};
  \ellips{0}{0}{2.449}{1}{0}
  \ellips{0}{0}{4.899}{2}{0}
  \punten[inkt]{(-0.47,0.65) (-1.41,-0.21) (0.32,-1.04) (1.43,-0.21) (0.91,-0.66) (3.16,0.58) (-0.34,0.50) (4.26,1.84) (-4.10,0.63) (1.07,-0.21) (-2.33,1.07) (-3.39,0.33) (2.60,-1.71) (-1.70,-1.55)}
  \node[below] at (0,-4.3) {gedraaid};
\end{tikzpicture}\quad
\begin{tikzpicture}[x=.4cm,y=.4cm,font=\scriptsize]
  \draw[->,grijs] (-4,0) -- (4.4,0) node[right,text=black] {$w_1$};
  \draw[->,grijs] (0,-4) -- (0,4.4) node[above,text=black] {$w_2$};
  \draw[accent,thick] (0,0) circle (1);
  \draw[accent,thick] (0,0) circle (2);
  \punten[inkt]{(-0.19,0.65) (-0.58,-0.21) (0.13,-1.04) (0.58,-0.21) (0.37,-0.66) (1.29,0.58) (-0.14,0.50) (1.74,1.84) (-1.67,0.63) (0.44,-0.21) (-0.95,1.07) (-1.38,0.33) (1.06,-1.71) (-0.70,-1.55)}
  \node[below] at (0,-4.3) {herschaald};
\end{tikzpicture}}

\voorbeeld{}{Bereken voor de wolk met $S = \km{5 & 2 \\ 2 & 2}$ de afstand van Mahalanobis van het punt met afwijking $\vec d = (3, 4)$, met de hoofdassen. Controleer met $S^{-1}$.\label{vb:h3-hoofdassen}}{%
\uitleg{In voorbeeld \ref{vb:h3-assen} vonden we $u_1 = 2\sqrt5$ en $u_2 = \sqrt5$. We delen elk kwadraat door zijn eigenwaarde.}
\begin{stappenbreed}
\dM^2 = \frac{20}{6} + \frac{5}{1} = \frac{25}{3}, \quad \dM \approx 2{,}89 & Met de hoofdassen. \\[2ex]
\dM^2 = \tfrac16\,(2 \cdot 9 - 4 \cdot 12 + 5 \cdot 16) = \tfrac{50}{6} = \tfrac{25}{3} & Met $S^{-1} = \frac16\km{2 & -2 \\ -2 & 5}$.
\end{stappenbreed}
Beide manieren geven hetzelfde. Langs de tweede hoofdas ligt het punt maar $2{,}24$ ver, maar daar is de standaardafwijking ook maar $\sqrt{1} = 1$: die richting telt het zwaarst.}

\begin{werkwijze}
Bereken de eigenwaarden en de genormeerde eigenvectoren van $S$. Bereken de nieuwe coördinaten $u_j = \vec e_j \cdot \vec d$ van de afwijking. Dan is $\dM^2 = \frac{u_1^2}{\lambda_1} + \frac{u_2^2}{\lambda_2}$. De ellips $\dM = k$ heeft halve assen $k\sqrt{\lambda_1}$ en $k\sqrt{\lambda_2}$ langs de eigenvectoren.
\end{werkwijze}

\begin{oefeningen}
\oef{h3-ellips}
\end{oefeningen}

\subsection{Beperkingen}

De afstand van Mahalanobis erft de beperkingen van de covariantiematrix. Ze bestaat alleen als $S$ inverteerbaar is: met te weinig individuen, of met een variabele die een lineaire combinatie is van de andere, lukt het niet. Ze veronderstelt dat de wolk ongeveer de vorm van een ellips heeft; bij een kromme of uit groepen bestaande wolk zegt ze weinig. En de uitschieters die je zoekt, zitten zelf in $S$: één ver punt maakt de wolk breder in zijn richting, en lijkt zo zelf minder ver. Daarom berekent men $S$ voor zo'n controle soms zonder het verdachte punt. Kijk dus, zoals altijd, ook naar de figuur.

\begin{oefeningen}
\oef{h3-beperking}
\end{oefeningen}

%========================================================================================
\needspace{16\baselineskip}
\subsection{Terug naar het probleem}

\voorbeeld{}{Terug naar het probleem. Van alle leerlingen van het zesde jaar kent de school de punten voor wiskunde ($W$) en fysica ($F$), in procent. Voor het hele jaar is het gemiddelde punt $(62, 58)$ en de covariantiematrix $S = \um{rr}{153 & 96 \\ 96 & 97}$. Lien haalt $83$ voor wiskunde en $80$ voor fysica, Ruben $74$ en $42$. De klassenraad wil weten welke leerling het meest afwijkt van wat gewoon is. Lien ligt het verst van het gemiddelde punt, maar is haar combinatie ook het meest ongewoon? Vergelijk de Euclidische, de gestandaardiseerde en de afstand van Mahalanobis.\label{vb:h3-terug}}{%
\uitleg{De afwijkingen tot $(62, 58)$ zijn $(21, 22)$ voor Lien en $(12, -16)$ voor Ruben. We gebruiken de hoofdassen: eerst de eigenwaarden en eigenvectoren van $S$.}
\begin{stappenbreed}
\text{Euclidisch: } \sqrt{21^2 + 22^2} \approx 30{,}4 \quad\text{en}\quad \sqrt{12^2 + 16^2} = 20 & Lien ligt het verst. \\[1ex]
\text{z-scores: } \sqrt{\tfrac{21^2}{153} + \tfrac{22^2}{97}} \approx 2{,}81 \quad\text{en}\quad \sqrt{\tfrac{12^2}{153} + \tfrac{16^2}{97}} \approx 1{,}89 & Nog altijd Lien. \\[2ex]
\lambda^2 - 250\lambda + 5625 = 0 \iff \lambda_1 = 225,\ \lambda_2 = 25 & $\det S = 14\,841 - 9216 = 5625$. \\[1ex]
\vec e_1 = \tfrac15(4, 3), \quad \vec e_2 = \tfrac15(-3, 4) & Uit $(S - \lambda I)\vec v = \vec 0$. \\[1ex]
\text{Lien: } u_1 = \tfrac{84 + 66}{5} = 30, \ u_2 = \tfrac{-63 + 88}{5} = 5, \ \dM^2 = \tfrac{900}{225} + \tfrac{25}{25} = 5 & $\dM \approx 2{,}24$. \\[1ex]
\text{Ruben: } u_1 = \tfrac{48 - 48}{5} = 0, \ u_2 = \tfrac{-36 - 64}{5} = -20, \ \dM^2 = 0 + \tfrac{400}{25} = 16 & $\dM = 4$.
\end{stappenbreed}
\naast{\begin{minipage}{9.2cm}
De afstand van Mahalanobis keert het oordeel om. Lien scoort hoog voor beide vakken: dat ligt in de lengterichting van de wolk, waar de standaardafwijking $\sqrt{225} = 15$ is, en ze ligt maar $2{,}24$ standaardafwijkingen van het centrum. Ruben scoort goed voor wiskunde maar zwak voor fysica, dwars op de wolk, waar de standaardafwijking maar 5 is: hij ligt 4 standaardafwijkingen ver, buiten de ellips $\dM = 3$. Zijn combinatie is de meest ongewone, en verdient de aandacht van de klassenraad. In de figuur zie je dat Lien binnen de ellips $\dM = 3$ ligt en Ruben erbuiten.
\end{minipage}}{\begin{tikzpicture}[x=.05cm,y=.05cm,font=\scriptsize]
  \diagram[$W$][$F$]{20}{10}{105}{20}{10}{100}
  \xas{20,40,...,100}\yas{20,40,...,100}
  \ellips{62}{58}{15}{5}{36.87}
  \ellips{62}{58}{30}{10}{36.87}
  \ellips{62}{58}{45}{15}{36.87}
  \punten[inkt]{(58,63) (49,45) (67,59) (55,53) (59,44) (43,44) (54,51) (65,63) (53,50) (62,62) (77,60) (59,52) (87,72) (71,57) (80,72) (52,49) (70,66) (54,50) (52,57) (67,57) (57,56) (53,43) (72,61) (82,74) (58,50) (43,49) (53,54) (55,44) (53,55) (82,75)}
  \fill[fel!80!black] (83,80) circle (2.4pt) node[above left,text=black] {Lien};
  \fill[fel!80!black] (74,42) circle (2.4pt) node[below right,text=black] {Ruben};
\end{tikzpicture}}}

%========================================================================================
\opdrachtenkop

\begin{oefeningen}
\oef{h3-o1}

\oef{h3-o2}

\needspace{10\baselineskip}
\oef{h3-o3}

\oef{h3-o4}

\needspace{10\baselineskip}
\oef{h3-o5}

\oef{h3-o6}

\needspace{12\baselineskip}
\oef{h3-o7}

\oef{h3-o8}

\needspace{12\baselineskip}
\oef{h3-o9}

\oef{h3-o10}

\needspace{12\baselineskip}
\oef{h3-o11}

\oef{h3-o12}

\needspace{12\baselineskip}
\oef{h3-o13}

\oef{h3-o14}

\needspace{12\baselineskip}
\oef{h3-o15}

\oef{h3-o16}

\needspace{12\baselineskip}
\oef{h3-o17}

\oef{h3-o18}

\needspace{12\baselineskip}
\oef{h3-o19}

\oef{h3-o20}

\needspace{12\baselineskip}
\oef{h3-o21}
\end{oefeningen}

%----------------------------------------------------------------------------------------
% KOPPEL: verbanden met andere leerstof
%----------------------------------------------------------------------------------------
\needspace{20\baselineskip}
\opdrachtsectie{Koppel}

\lkr{De eerste oefening steunt op Kegelsneden (de vergelijking van een ellips in haar assenstelsel, de brandpunten), de tweede op Matrices (hoofdstuk 4: overgangsmatrices en de evenwichtstoestand), de derde op Goniometrische functies (de verdubbelingsformules en $a\cos x + b\sin x$ als één sinusoïde), de vierde op de module Eigenwaarden en eigenvectoren van Complementair (LPD 30, hoofdstuk 3: machten van een diagonaliseerbare matrix); geef die alleen als de leerlingen ze zagen.}

\begin{oefeningen}
\oef{h3-koppel-1}

\oef{h3-koppel-2}

\oef{h3-koppel-3}

\oef{h3-koppel-4}
\end{oefeningen}

%----------------------------------------------------------------------------------------
\opdrachtsectie{Test jezelf}

\begin{oefeningen}
\oef{h3-t1}

\oef{h3-t2}

\oef{h3-t3}

\oef{h3-t4}

\oef{h3-t5}

\oef{h3-t6}

\oef{h3-t7}
\end{oefeningen}

\herhaalschema{%
  \fv{oef:h3-t1}{vb:h3-mahalanobis}{oef:h3-mahalanobis}%
  \fv{oef:h3-t2}{vb:h3-eigen}{oef:h3-eigen}%
  \fv{oef:h3-t3}{vb:h3-diagonaliseren}{oef:h3-diagonaliseren}%
  \fv{oef:h3-t4}{vb:h3-assen}{oef:h3-assen}%
  \fv{oef:h3-t5}{vb:h3-hoofdassen}{oef:h3-ellips}%
  \fv{oef:h3-t6}{vb:h3-hoofdassen}{oef:h3-mahalanobis}%
  \fv{oef:h3-t7}{vb:h3-terug}{oef:h3-beperking}}

\ifzelfstudie
\section*{Maak je eigen samenvatting}
Sluit dit boek en schrijf in je schrift: waarom de Euclidische afstand niet volstaat (eenheden en samenhang), de definitie van de afstand van Mahalanobis en de formule bij twee variabelen, de twee eigenschappen met hun bewijs, de ellipsen van gelijke afstand, wat een eigenwaarde en een eigenvector zijn, de karakteristieke vergelijking, de drie eigenschappen van symmetrische matrices en covariantiematrices met hun bewijs, het diagonaliseren $S = PDP\tr$, waarom de hoofdcomponenten ongecorreleerd zijn, waarom de eerste hoofdas de grootste spreiding geeft, de afstand van Mahalanobis in de hoofdassen en de halve assen van de ellips, en de beperkingen. Vergelijk daarna met de kaders.
\fi
